\chapter{Synthesis of the Results} \label{ch:synthesis}

While Chapter \ref{ch:results} provided a detailed quantitative analysis of the simulation data, this chapter elevates the findings to a broader meta-level. The purpose of this synthesis is to directly answer the core research question, critically evaluate the strengths and limitations of the applied EPEC methodology, discuss the transferability of the findings to real-world power systems, and reflect on the lessons learned during the modeling process.

\section{Answers to the Research Question} \label{sec:answers_rq}

The overarching research question of this thesis explored what insights into strategic investor behavior can be gained by shifting from a traditional central-planner cost-minimization approach to a decentralized, game-theoretic EPEC framework. Based on the quantitative results, this question can be answered clearly: Traditional cost-minimization models implicitly assume perfect system-wide coordination and unselfish asset deployment. In contrast, the EPEC framework uncovered that strategic, profit-maximizing behavior fundamentally alters both the operational logic and the physical dimensioning of BESS.

Regarding the \textbf{siting (location)} of BESS, the results demonstrate that strategic investors distribute their assets across the entire grid. Rather than clustering exclusively at nodes with the highest renewable generation (e.g., wind at Node 1 or PV at Node 3), the capacities were allocated almost evenly across all available nodes. This spatial distribution is the direct result of fierce oligopolistic competition for highly lucrative market opportunities. Because participation in the system-wide markets -- especially the reserve market -- is overwhelmingly profitable, every single node represents a valuable access point to generate revenue. Consequently, to maximize their individual market shares and overall profits, the investors aggressively compete for, and fully exhaust, the available grid connection limits at every node across the network.



Regarding the \textbf{sizing (capacity)} of BESS, the framework revealed a fundamental divergence between power capacity (MW) and energy capacity (MWh). The optimal power capacity is primarily dictated by the competitive environment. Under the continuous Cournot Inverse Demand mechanism, the investors formed a balanced oligopoly, sharing the power capacity provision equally to keep endogenous reserve prices high and avoid market cannibalization.

In contrast, the optimal energy capacity (the E/P ratio) is heavily dictated by the individual economic constraints of the investors -- specifically, their cost of capital (WACC) -- and the overwhelming profitability of the reserve market. The model proved that "patient" investors with low capital costs (e.g., 8\%) deploy massive energy reservoirs (E/P ratio $> 21$) to compensate for a severe power bottleneck. Because the expensive inverter capacity (MW) is almost entirely blocked for reserve provision, investors oversize their cheap energy capacity (MWh) to still extract marginal revenues out of the spot market with the remaining fractional power.

Ultimately, the inclusion of both spot and reserve markets into the modeling framework uncovered a severe cross-market spillover effect. The results answer the research question by highlighting that optimal BESS allocation in a multi-service environment is not solely a function of geographic arbitrage. Instead, it is governed by an opportunity-cost logic: the strategic lock-up of power capacity for highly lucrative reserve markets forces the storage units into inefficient, slow-cycling, and occasionally loss-making operations on the spot market. Consequently, the true value of a BESS location and size is fundamentally shaped by the specific regulatory and economic design of the ancillary service markets it participates in.

\section{Strengths and Limitations of the EPEC Methodology} \label{sec:strengths_limitations}

The application of the EPEC framework to model the strategic allocation of BESS in a multi-service market proved to be both a powerful analytical tool and a significant mathematical challenge. By critically evaluating the modeling process, several distinct methodological strengths and necessary computational compromises can be identified.

\subsection{Strengths of the Applied Framework}
The primary strength of the EPEC methodology lies in its ability to simulate true oligopolistic competition through endogenous price formation. Unlike standard cost-minimization models where market prices are either exogenous inputs or remain independent of a single actor's investment scale, the EPEC actively internalizes the market's price response to newly installed BESS capacities. 

During the development, the implementation of the \textit{Residual Demand} formulation proved to be a critical methodological strength. By allowing investors to strictly optimize against the pre-calculated residual load of their rivals, the framework successfully eliminated the "cobweb" oscillations typically observed in iterative equilibrium models. Furthermore, the integration of a continuous \textit{Cournot Inverse Demand Curve} for the aFRR market allowed for a realistic, stable representation of cartel-like profit maximization. This continuous formulation successfully prevented the non-linear solver from exploiting mathematical KKT-loopholes (the "price-dictator" exploit), ensuring that the resulting equilibrium is rooted in sound economic theory rather than mathematical artifacts.

\subsection{Limitations and Mathematical Compromises}
Despite these strengths, the non-linear and non-convex nature of coupled KKT conditions required several structural compromises to maintain computational tractability within the applied interior-point solver (Ipopt):

\begin{itemize}
    \item \textbf{Continuous Approximation of Pay-as-Bid Markets:} Real-world ancillary service markets, such as the Austrian APG aFRR market, utilize a Pay-as-Bid clearing mechanism and require discrete minimum bid sizes (e.g., integer MW steps) \cite{apg_secondary_control}. Implementing a true Pay-as-Bid auction would have transformed the model into a Mixed-Integer Non-Linear Program with Equilibrium Constraints (MI-NLP). Given the state-of-the-art of optimization solvers, solving an MI-NLP of this scale within an iterative diagonalization loop is computationally intractable. Consequently, the continuous Cournot formulation was a mathematically necessary macroeconomic approximation.
    \item \textbf{KKT Relaxation and Numerical Noise:} To navigate the severe mathematical singularities that occurred during extreme peak-price and zero-price (curtailment) events, the strict complementarity conditions of the lower level had to be significantly relaxed ($\epsilon = 10.0$). While this provided the solver with the necessary gradient tube to ensure algorithmic convergence, it introduced microscopic numerical noise (e.g., fractional generator dispatch despite 0 €/MWh spot prices).
    \item \textbf{Linearized Battery Degradation:} Actual lithium-ion battery degradation is a highly complex, non-linear process heavily dependent on the Depth of Discharge (DoD) and specific cycle profiles, typically evaluated using Rainflow-Counting algorithms \cite{shin_optimal_2020}. Because such algorithmic evaluations are non-differentiable, they cannot be embedded into an NLP solver's constraints. Therefore, degradation had to be abstracted as a constant linear marginal cost per MWh throughput, slightly simplifying the operational reality of BESS.
    \item \textbf{Perfect Foresight (Deterministic Approach):} The model assumes perfect foresight regarding renewable generation and load profiles. Introducing stochastic optimization to account for forecast errors and actual balancing energy activation would have caused an exponential explosion in computation time, rendering the EPEC unsolvable on standard hardware.
    \item \textbf{Exclusion of Grid Tariffs and Network Fees:} In real-world electricity markets, BESS assets are often subjected to grid fees, taxes, and levies during their charging and discharging cycles. Because a BESS acts as both a consumer and a generator, these fees essentially act as a financial penalty on energy throughput, significantly shrinking the available arbitrage margins \cite{hameed_business-oriented_2021}. The current framework assumes a frictionless market without these grid fees, which leads to a slight overestimation of the profitability of pure spot market arbitrage.
    \item \textbf{Omission of Operations and Maintenance (O\&M) Costs:} While the model rigorously accounts for annualized capital expenditures (CAPEX) and cyclic battery degradation, fixed annual O\&M costs (typically estimated at around 1\% of CAPEX \cite{shin_optimal_2020}) were deliberately excluded from the objective function.


\end{itemize}

\section{Upscaling and Transferability of Results} \label{sec:upscaling}

The EPEC modeling framework developed in this thesis, calibrated on a 5-bus network over a representative 24-hour period, serves as a rigorous mathematical proof-of-concept. It successfully demonstrates the intricate strategic dynamics of BESS investments in multi-service markets. However, rather than serving as a real-time operational tool, the primary value of this framework lies in its potential for advanced academic research and macroeconomic forecasting. To utilize this framework for comprehensive forecasts of future energy systems, long-term scenario planning, and data-driven policy consulting, several dimensions of upscaling must be addressed in future research.

\subsection{Temporal and Spatial Upscaling}
Currently, the model's temporal horizon is restricted to a single, highly volatile representative day. While this accurately captures diurnal arbitrage and physical battery cycling, it inherently neglects seasonal variations. Renewable generation patterns -- such as prolonged dunkelflauten versus extended summer photovoltaic peaks -- fundamentally alter the long-term profitability and state-of-charge management of BESS. Future applications of this model should upscale the temporal resolution by clustering historical data into seasonal representative days (e.g., distinct profiles for spring, summer, autumn, and winter) or utilizing a rolling-horizon approach to approximate an 8760-hour year. 

Spatially, the transferability of the model to large-scale, real-world power systems (such as the IEEE 39-bus system or the complete Austrian high-voltage transmission grid) is currently limited by the computational tractability of EPECs. As the number of nodes, lines, and strategic actors increases, the dimensionality of the coupled KKT conditions grows exponentially. Upscaling this framework to national grid levels will likely require the integration of advanced decomposition algorithms alongside the currently implemented diagonalization heuristic.

\subsection{Transferability to European Market Designs}
A critical aspect of this framework's transferability is its foundational market design. The implemented DC-OPF market clearing natively generates Locational Marginal Prices (LMPs), reflecting a \textit{Nodal Pricing} system. While this is the standard market design in various U.S. markets, the European electricity market currently operates on a zonal pricing model ("copper plate" assumption), where internal grid congestion is resolved post-market via costly redispatch measures. 

Therefore, the results of this thesis offer a powerful, forward-looking perspective. They demonstrate how strategic BESS investments would organically distribute themselves across the grid to relieve structural bottlenecks \textit{if} a Nodal Pricing system were implemented. The framework can be highly valuable for the debate of transitioning towards smaller bidding zones or full nodal pricing, as it quantifies the precise economic incentives required to stimulate grid-supportive BESS allocation without relying on central-planner redispatch interventions.

\subsection{Algorithmic Evolution towards MIPECs}
Finally, the transferability of the operational results to real-world ancillary service markets is bounded by the continuous relaxation of the bidding variables. Real reserve markets demand discrete bid sizes (e.g., 1 MW increments). While the continuous \textit{Cournot Inverse Demand} approximation utilized in this thesis is currently the most robust macroeconomic workaround, the ultimate goal for future research is the formulation of a Mixed-Integer Program with Equilibrium Constraints (MIPEC). Once commercial MI-NLP (Mixed-Integer Non-Linear Programming) solvers reach a maturity level capable of handling massive, coupled complementarity conditions, the framework can be upgraded to reflect exact, discrete market auction rules without succumbing to the mathematical exploits currently observed in non-linear environments.

\subsection{Transferability to Co-located Hybrid Configurations}
The current framework exclusively simulates standalone BESS, where each unit interacts directly with the high-voltage transmission grid. However, in modern real-world project development, BESS units are increasingly co-located with renewable energy generators (such as solar PV and wind) in hybrid park configurations, sharing a single grid connection point. Transferring the EPEC framework to model such topologies represents a highly relevant future upscaling step. In a hybrid setup, the BESS can charge directly from local renewable surplus during high-yield periods without utilizing the public grid, thereby entirely avoiding transmission grid fees. Integrating these behind-the-meter dynamics would add a critical layer of strategic complexity, as investors would need to co-optimize grid-fee-free local charging against system-wide market arbitrage and reserve provision.

\section{Lessons Learned from EPEC Modeling} \label{sec:lessons_learned}

The development and deployment of the EPEC framework for this thesis provided profound insights not only into BESS allocation but also into the fundamental mechanics of advanced energy system modeling. Reflecting on the iterative development process, several critical lessons learned emerged that can serve as best practices for future research in this highly complex field.

\subsection{Algorithmic Limitations in Capturing Economic Rationality}
A primary realization was that non-linear optimization solvers (such as Ipopt) are strictly mathematical engines that do not inherently account for economic rationale. If a mathematical loophole exists, the solver will exploit it to maximize the objective function, often leading to market-distorting paradoxes. This was profoundly evident in two instances during model development: 
First, the Zero-Price exploit in the spot market, where MPEC investors simulated microscopic, fictitious renewable curtailment using KKT slack tolerances solely to force the nodal market price artificially to 0 €/MWh, allowing them to charge their batteries for free. 
Second, the Price-Dictator phenomenon in the reserve market, where investors supplied near-zero reserve capacity but utilized the KKT complementarity slack to push the clearing price to the absolute penalty cap (3,000 €/MW). 
These findings demonstrate that mathematical formulations in energy modeling should be closely guided by underlying economic principles. Replacing rigid KKT market clearing constraints with continuous economic mechanisms -- such as the implemented Cournot Inverse Demand curve and penalty-driven soft constraints -- is not a dilution of the model, but a mandatory requirement to ensure rational market behavior.

\subsection{Challenges in Computing Simultaneous Market Equilibria}
Modeling multiple strategic actors in a highly constrained power grid inherently invites the "cobweb" oscillations. Early model iterations revealed that when investors optimize based on the static, fixed decisions of their rivals, they continually evade each other in consecutive iterations, preventing the algorithm from converging to a Nash Equilibrium. The implementation of the \textit{Residual Demand} concept -- where an investor optimizes against the already shifted residual load profile of the entire system rather than the isolated actions of peers -- was a fundamental breakthrough. It proved that in EPEC modeling, agents must react to the \textit{net market state} rather than individual opponent variables to achieve algorithmic stability.

Finally, the development process highlighted that insisting on absolute mathematical purity often leads to computationally intractable models. In highly volatile power systems characterized by steep peak-load prices and zero-price renewable events, strict KKT complementarity conditions create severe mathematical singularities that crash interior-point solvers. Embracing pragmatic engineering solutions -- such as relaxing the KKT conditions (e.g., $\epsilon = 10.0$), utilizing algorithmic damping factors ($\alpha = 0.25$) to suppress oscillation, and applying ex-post filters to remove microscopic numerical noise -- proved absolutely essential. These compromises are not methodological flaws, but rather the necessary state-of-the-art tools required to make the theoretical elegance of EPECs practically applicable.

